%% kapa: exam=1 round=31 year=2020 period=2 form=A subject=appraisal_law subject_ko=감정평가관계법규 seq=20 number=20
%% source: Q-Net 감정평가사 기출문제 (cst00309) artlSeq=5207219 file=extracted/1008/5207219/A형/A형/2020년도 제31회 감정평가사 1차 시험_2교시_A형.pdf
%% conversion: kordoc 4.24.51 -> tools/gen1.py (auto); review=proofread-sonnet; flags=none
\begin{question}{20}
국유재산법령상 행정재산에 관한 설명으로 옳지 않은 것은?
\begin{choices}
  \item 행정재산은 사유재산과 교환할 수 없다.
  \item 행정재산을 경쟁입찰의 방법으로 사용허가하는 경우 1개 이상의 유효한 입찰이 있으면 최고가격으로 응찰한 자를 낙찰자로 한다.
  \item 두 번에 걸쳐 유효한 입찰이 성립되지 아니한 행정재산의 경우 수의의 방법으로 사용허가를 받을 자를 결정할 수 있다.
  \item 중앙관서의 장이 행정재산의 사용허가를 철회하려는 경우에는 청문을 하여야 한다.
  \item 중앙관서의 장은 행정재산으로 사용하기로 결정한 날부터 5년이 지난 날까지 행정재산으로 사용되지 아니한 행정재산은 지체 없이 그 용도를 폐지하여야 한다.
\end{choices}
\end{question}
\answer{1}
